\subsection{THETA3-C37. Simultaneous rank-one differences of a positive theta section}
\label{C37}

This result supplies a positive analytic rank theorem for several
incidence factors at the same rational prime.  It allows the actual
torsion representatives to be unbounded.  Its application to a particular
finite transform still requires that transform to have the ordered
operator form stated below and to have one fixed nonzero baseline theta
section.  The graph-subgroup argument was developed jointly with the
independent review of A15; none of the phase factors is discarded.

\paragraph{The analytic and sampling hypotheses.}
Let $T=V/\Lambda$ be a complex torus of positive dimension, with a
positive definite Hermitian form $H$, linear in its first argument.
Write $E=\operatorname{Im}H$ and assume $E(\Lambda,\Lambda)\subset\Z$.
Let $L$ be the positive Appell--Humbert line with semicharacter
$\alpha$, and let $F\ne0$ be an entire theta function with the
normalization
\[
 \alpha(\lambda)U_\lambda F=F,
 \qquad
 U_yF(z)=F(z+y)
       \exp\bigl(-\pi H(z,y)-\tfrac\pi2H(y,y)\bigr).
 \tag{C37a}
\]
Fix an integer $f\ge1$.  For each sufficiently large index $q$, and
each $i=1,\ldots,f$, let $S_{i,q}$ be a finite nonempty sample set,
with a map $s\mapsto t_{i,q}(s)\in T$ and uniform probability measure
$\mu_{i,q}$ on these points, counting multiplicities.  Suppose that
there are numbers $\epsilon_q\to0$ with the following two properties.
For every character $\gamma\in\Lambda^*=\Hom(\Lambda,\Z)$,
\[
 \widehat\mu_{i,q}(\gamma)=1
 \quad\hbox{or}\quad
 |\widehat\mu_{i,q}(\gamma)|\le\epsilon_q;
 \tag{C37b}
\]
the first alternative occurs precisely when the character is trivial
on the underlying finite torsion group.  For each fixed nonzero
$\gamma$, that first alternative does not occur for all sufficiently
large $q$.  The latter condition implies
$\mu_{i,q}\to m_T$, the probability Haar measure, by approximation
with trigonometric polynomials.

Choose an arbitrary lift $Y_{i,q}(s)\in V$ of every sample point and
an arbitrary scalar $\eta_{i,q}(s)$ of modulus one.  The lifts need
not lie in a bounded set.  Consider the first-derivative columns
\[
 w_q(s_1,\ldots,s_f)=
 d_z\left[
   (\eta_{f,q}(s_f)U_{Y_{f,q}(s_f)}-1)\cdots
   (\eta_{1,q}(s_1)U_{Y_{1,q}(s_1)}-1)F
       \right]_{z=0}\in V^*.
 \tag{C37c}
\]
Each phase here depends on its own factor.  Bilinear phases produced
by reducing the lifts will be retained in the proof.

\paragraph{The simultaneous spanning theorem.}
For every sufficiently large $q$, the columns (C37c) span $V^*$.
More precisely, choose points $t_2,\ldots,t_f\in V$ such that
\[
 \sum_{i\in J}t_i-\sum_{i\in K}t_i\notin\Lambda\otimes\Q
 \quad\text{whenever }J\ne K\subset\{2,\ldots,f\}.
 \tag{C37d}
\]
Such choices exist: the prohibited set is a countable union of
proper real affine subspaces in $V^{f-1}$.  Equidistribution permits
choices $s_{i,q}\in S_{i,q}$, for $i\ge2$, whose images converge to
$t_i$ modulo $\Lambda$.  Then there is a constant $c>0$ such that,
for all sufficiently large $q$, the matrix $J_q$ formed from the
columns $w_q(s_1,s_{2,q},\ldots,s_{f,q})$ satisfies
\[
 J_qJ_q^*\ \ge\ c\,|S_{1,q}|\,\operatorname{Id}_{V^*}.
 \tag{C37e}
\]
Here the Hermitian norm on $V^*$ is fixed.  The conclusion is
uniform in all choices of lifts and phases in (C37c).  The proof
gives an existence threshold and a positive constant, rather than
an explicit numerical threshold for these several-factor data.
For $f=1$, A15 gives the explicit constant $\delta_F/2$ and its
Fej\'er threshold.

We prove (C37e) by contradiction, recording the two geometric
ingredients in full.

\paragraph{Reduction of the lifts and its exact extra characters.}
For $\lambda\in\Lambda$, direct use of (C37a) gives
\[
 U_{y+\lambda}F
  =\alpha(\lambda)^{-1}e^{\pi iE(y,\lambda)}U_yF.
 \tag{C37f}
\]
The multiplier is independent of $z$.  Thus for every fixed linear
operator $P$ on these entire functions and tangent vector $u$,
$d(PU_yF)_0(u)$ is a smooth section of the same theta line in the
parameter $y$; its norm is well defined on $T$.  In particular,
it is not being treated as a holomorphic scalar function of $y$.

For a translation on the left the exact identity is
\[
 U_{t+\lambda}U_yF
   =\alpha(\lambda)^{-1}e^{\pi iE(t,\lambda)}
       e^{-2\pi iE(\lambda,y)}U_tU_yF.
 \tag{C37g}
\]
Indeed $U_xU_y=e^{-\pi iE(x,y)}U_{x+y}$, and (C37f) applied
to $x+y$ proves (C37g).  The last factor is a character of $T$,
because $E(\lambda,\Lambda)\subset\Z$.  This character, which
can vary with $q$, is essential.

For the chosen samples in factors $i\ge2$, choose bounded lifts
$t_{i,q}\to t_i$ and write their actual lifts as
$Y_{i,q}=t_{i,q}+\lambda_{i,q}$, $\lambda_{i,q}\in\Lambda$.
Put $t_{J,q}=\sum_{i\in J}t_{i,q}$.  Expanding the ordered product
of the last $f-1$ differences and applying (C37g) gives, for every
tangent vector $u$,
\[
 d\left[
  (\eta_fU_{Y_f}-1)\cdots(\eta_2U_{Y_2}-1)U_yF
       \right]_0(u)
   =\sum_{J\subset\{2,\ldots,f\}}
         c_{J,q}\prod_{i\in J}\chi_{i,q}(y)\,
               d(U_{t_{J,q}}U_yF)_0(u),
 \tag{C37h}
\]
where $|c_{J,q}|=1$ and
\[
       \chi_{i,q}(y)=e^{-2\pi iE(\lambda_{i,q},y)}.
 \tag{C37i}
\]
All signs, the factors $\eta_i$, the Heisenberg cocycles, and the
constant multipliers in (C37g) are included in $c_{J,q}$.
There is exactly one summand per subset, so none of these
coefficients vanishes.

For a first-factor sample, (C37f) removes its lattice displacement
before the fixed left operator is applied.  Consequently its
pairing with $u$ has the form
\[
       \zeta_{s,q}A_q(t_{1,q}(s),\chi_{2,q}(t_{1,q}(s)),\ldots)
                       -A_q(0,1,\ldots,1),
       \qquad |\zeta_{s,q}|=1,
 \tag{C37j}
\]
where $A_q(y,z_2,\ldots,z_f)$ is the right side of (C37h) with
$\chi_{i,q}(y)$ replaced by $z_i$.  This is a section of the
pullback of $L$ to $T\times(S^1)^{f-1}$.  Taking absolute values
in (C37j) eliminates only the first scalar $\zeta_{s,q}$; it does
not eliminate the characters (C37i).

\paragraph{Haar limits on graph subgroups.}
Push $\mu_{1,q}$ forward under
\[
       y\longmapsto(y,\chi_{2,q}(y),\ldots,\chi_{f,q}(y))
          \in T\times(S^1)^{f-1},
\]
and denote the resulting probability measure by $\nu_q$.
Every subsequential weak limit of these measures is Haar measure
on a closed subgroup $G\subset T\times(S^1)^{f-1}$ whose
projection to $T$ is surjective.  Here is an elementary proof.
For a character $(\gamma,k_2,\ldots,k_f)$ of this product, its
integral against $\nu_q$ is a Fourier coefficient of $\mu_{1,q}$
at the possibly moving character
$\gamma+\sum_i k_i\chi_{i,q}$.  By (C37b), every limit of these
coefficients is either $0$ or $1$.  If the integral of a
unit-modulus character against a probability measure equals $1$,
that character equals $1$ almost everywhere.  The characters with
integral $1$ therefore form a subgroup $D$ of the character group.
The limiting measure and Haar measure on the annihilator of $D$
have identical Fourier coefficients, hence coincide, by density
of trigonometric polynomials.  This proves the subgroup assertion.
For characters with $k_i=0$, the fixed-character part of (C37b)
shows that the marginal is $m_T$.  The projection of $G$ is
closed, and therefore is all of $T$.

For later use, any such projection $G\to T$ has a section after
a finite identity-linear covering of $T$.  Indeed $G$ is a
compact abelian Lie group.  Its identity component $G^0$ still
surjects onto $T$: the quotient by the image of $G^0$ is both
connected and finite.  The surjection of the rational fundamental
groups of these tori has a rational linear right inverse.
Clearing its denominators gives a homomorphism
\[
 j:V/N\Lambda\longrightarrow G^0,
       \operatorname{pr}_Tj(y)=y\pmod\Lambda
 \tag{C37k}
\]
for some positive integer $N$.  Each circle coordinate of this
homomorphism has the form $e^{2\pi i\ell_i(y)}$, where
$\ell_i\in\Lambda^*\otimes\Q$.  This argument does not assume
that $G$ is connected or that its projection has a section
without a finite covering.

\paragraph{Nonvanishing on every such subgroup.}
Let $u\ne0$, let all $c_J$ have modulus one, and set
\[
 A(y,z)=\sum_J c_J\prod_{i\in J}z_i\,
                       d(U_{t_J}U_yF)_0(u),
       \qquad t_J=\sum_{i\in J}t_i.
 \tag{C37l}
\]
For every closed subgroup $G$ projecting onto $T$, the norm of
this section is not constant on $G$.

First we prove that its restriction along (C37k) is not the zero
section.  Since the rational alternating form $E$ is
nondegenerate, for every $\ell_i\in\Lambda^*\otimes\Q$ there
is $s_i\in\Lambda\otimes\Q$ with $\ell_i(y)=E(s_i,y)$.
Put $s_J=\sum_{i\in J}s_i$.  Differentiating (C37a) gives
\[
 d(U_tU_yF)_0(u)=
 e^{-\pi H(y,y)/2-\pi H(t,t)/2-\pi H(t,y)}
 \bigl(dF_{y+t}(u)-\pi H(u,y+t)F(y+t)\bigr).
 \tag{C37m}
\]
The equality
$2\pi iE(s,y)=\pi H(s,y)-\pi H(y,s)$ shows that, after
pullback along $j$ and multiplication by $e^{\pi H(y,y)/2}$,
the summands have antiholomorphic exponential slopes
$-\pi H(t_J-s_J,y)$.  These slopes are distinct: equality for
$J\ne K$ would contradict (C37d).

For completeness, the required exponential independence can be
proved by differential operators.  Choose a direction $v\in V$
such that $H(u,v)\ne0$ and all numbers
\[
           \rho_J=-\pi H(t_J-s_J,v)
\]
are distinct; only finitely many proper complex hyperplanes
are excluded.  Write $D=\partial_{\bar v}$.  Each summand is
an antiholomorphic exponential times a polynomial of degree at
most one in the antiholomorphic variables, with holomorphic
coefficients.  Thus $(D-\rho_J)^2$ kills that summand.
Its image under $D-\rho_J$ is the nonzero real-analytic function
\[
 -\pi H(u,v)c_J
 e^{-\pi H(t_J,t_J)/2-\pi H(y,s_J)}
 e^{-\pi H(t_J-s_J,y)}F(y+t_J).
\]
If the sum vanished identically, applying
$(D-\rho_J)\prod_{K\ne J}(D-\rho_K)^2$ would leave this
nonzero function multiplied by
$\prod_{K\ne J}(\rho_J-\rho_K)^2$, a contradiction.
This proves nonvanishing without a genericity assumption on $F$.

If the norm in (C37l) were constant on $G$, its pullback along
$j$ would have constant norm.  A zero constant contradicts the
nonvanishing just proved.  A positive constant supplies a
nowhere-zero smooth section of the pullback of $L$ to
$V/N\Lambda$, trivializing that line topologically.  This is
impossible: its positive curvature has nonzero top intersection,
whereas a topologically trivial line has zero first Chern class.
This proves the assertion about the norm on $G$.

\paragraph{Completion of the uniform rank proof.}
If (C37e) failed for arbitrarily large $q$ with no positive lower
bound, there would be a sequence of unit tangent vectors $u_q$
for which
\[
  \frac1{|S_{1,q}|}\sum_{s\in S_{1,q}}
       |\langle u_q,w_q(s,s_{2,q},\ldots,s_{f,q})\rangle|^2
                  \longrightarrow0.
 \tag{C37n}
\]
Pass to subsequences so that $u_q\to u$ is a unit vector,
every $c_{J,q}\to c_J$, and $\nu_q\to m_G$ as above.
The bounded translations already satisfy $t_{J,q}\to t_J$.
In local unitary frames the corresponding sections $A_q$ converge
uniformly to (C37l); their transition functions are all those of
the same line $\operatorname{pr}_T^*L$, by (C37f).  Thus the
squared norm-difference functions converge uniformly as well.
The inequality $|\zeta a-b|\ge\bigl||a|-|b|\bigr|$ and
(C37j)--(C37n) imply
\[
       \int_G\bigl(|A(y,z)|-|A(0,1,\ldots,1)|\bigr)^2\,dm_G=0.
 \tag{C37o}
\]
Continuity implies that the integrand vanishes everywhere on
$G$, contrary to the preceding paragraph.  This proves (C37e),
including uniformity over arbitrary lifts and phases.  A
minimum-norm right inverse of $J_q$ has coefficient-sum bound
$c^{-1/2}\|w\|$ for a desired first derivative $w$, by
Cauchy--Schwarz and (C37e).

\paragraph{The exact CM rank-one samples.}
Here are sufficient arithmetic hypotheses for (C37b).  Let
$\Lambda$ be a full lattice over an order in a CM field $E_0$,
with maximal order and perfect trace pairing at the primes under
consideration.  For a prime ideal $\mathfrak p$, the ideal-torsion
group is $\mathfrak p^{-1}\Lambda/\Lambda$.  Define the action
on the rational trace dual by $(a\ell)(x)=\ell(ax)$.  Its exact
annihilator in $\Lambda^*$ is
\[
 \{\ell:\ell(\mathfrak p^{-1}\Lambda)\subset\Z\}
                        =\mathfrak p\Lambda^*.
 \tag{C37p}
\]
This follows locally by dualizing the invertible ideal
$\mathfrak p^{-1}$, and away from $\mathfrak p$ both lattices
are unchanged.  Under an identification using a Hermitian
polarization, conjugate ideals may appear instead; (C37p) uses
the stated contragredient trace-dual convention.  A fixed nonzero
dual element lies in $\mathfrak p\Lambda^*$ for only finitely
many primes, as is seen from any nonzero coordinate in a finite
generating set after clearing denominators.

Assume the actual residue frames identify this torsion group with
$\operatorname{Mat}_{b\times a}(\mathbb F_q)$ and its dual
pairing with $\psi(\operatorname{Tr}(A^tX))$.  For the uniform
sample of nonzero rank-one matrices, a direct sum over
$A=uv^t$, $u,v\ne0$, gives
\[
 \widehat\mu(X)=
 \frac{q^{a+b-r}-q^a-q^b+1}{(q^a-1)(q^b-1)},
                  \qquad r=\operatorname{rank}X.
 \tag{C37q}
\]
Every rank-one matrix has $q-1$ such presentations.  If $r=0$
this coefficient is $1$; if $r\ge1$ its absolute value is at
most $1/(q-1)$.  The latter follows by counting the
$q^{b-r}-1$ nonzero vectors in the left kernel and using that
the inner additive-character sum is $q^a-1$ there and $-1$
elsewhere.  Thus (C37b) holds with $\epsilon_q=1/(q-1)$,
including for the moving characters in the graph proof.
Equation (C37p) proves the required fixed-character condition.
Different selected ideals above the same rational prime are
allowed; the proof never chooses the prime for one factor after
fixing a prime for another factor.

\paragraph{Scope for the literal finite transform.}
If the finite Fourier transform sends each factor character to
$T_y=\alpha(y)U_y$ with $y$ in a common half lattice $B$ and
$E(B,B)\subset\Z$, then the semicharacter relation and (C37a)
give $T_xT_y=T_{x+y}$.  Hence a product of character differences
has exactly the ordered form (C37c).  The rank-one incidence
identity (A15j) then converts the columns back to the permitted
incidence span; it does not enlarge that span.  Source and target
residue frames and the different factor must be the actual
tensor-compatible ones, as in (A15p).

This conditional application retains two exact checks.  The
remaining local factors must give a fixed nonzero positive theta
section $F$, and the actual transform must have the asserted
ordered character-operator realization.  A nonzero raw finite
array is not a proof of the first check.  If the output is viewed
on $V/N\Lambda$, the coordinate is still the $z$ of (C37a).
Identifying that cover with $V/\Lambda$ by $z=Nw$ multiplies
first derivatives by $N$, exactly as in A15.  No such factor,
canonical theta period, arithmetic normalization, or later
cohomological comparison is absorbed in the analytic theorem.
